Il·luminem una superfície de coure amb llum de diverses freqüències i quan s'alliberen electrons del metall, en mesurem l'energia cinètica. Amb les dades obtingudes de l'experiment dibuixem la gràfica següent:

\figura[0.45]{fig1.png}

\begin{apartat}{1}
Expliqueu breument què és el \textit{llindar de freqüència} de l'efecte fotoelèctric i calculeu quin valor té en aquest cas.
\end{apartat}

\begin{apartat}{1}
Calculeu el valor de la constant de Planck i la velocitat que assoleixen els electrons emesos quan la longitud d'ona de la llum incident és $1,2\times 10^{-7}\ \mathrm{m}$.
\end{apartat}

\dades{%
  $c = 3,00\times 10^{8}\ \mathrm{m\,s^{-1}}$\\
  $m_\text{electró} = 9,11\times 10^{-31}\ \mathrm{kg}$\\
  $1\ \mathrm{eV} = 1,60\times 10^{-19}\ \mathrm{J}$}
